% !TeX root = ../main.tex
% !TeX encoding = UTF-8
% !TeX spellcheck = none
\subsection{消融实验}\label{sec:exp-ablation}

为评估各修正通路的贡献，我们在成都T5数据集上构建三种消融变体：去除自身专家的A1、去除有向时滞专家的A2，以及去除BCMS的A3。各变体均独立重新训练。对于A1与A2，被移除专家的门控权重置零，剩余专家的权重归一化为1。测试沿用第\ref{sec:exp-baselines-metrics}节的正目标值掩码与AVG12统计口径，以流入、流出指标的等权平均衡量整体预测性能；分通道结果见补充材料S4。

表~\ref{tab:serr-expert-ablation}显示，完整CMSRN在联合MAE和联合RMSE上均取得最低误差，三种消融变体则呈现不同的性能变化。去除自身专家后，联合MAE增加0.386\%，为三种变体中的最大增幅，同时联合RMSE增加0.301\%。这一结果说明，节点自身时序演化为基础预测提供了重要的修正依据，对降低整体平均误差具有较大贡献。

去除有向时滞专家后，联合MAE增加0.044\%，联合RMSE增加0.362\%，后者为三种变体中的最大增幅。考虑到RMSE对较大误差更敏感，这一差异表明，时滞关联的价值更多体现在对较大预测偏差的修正上。自身专家与时滞专家在两项指标上呈现不同的贡献特点，支持分别组织节点自身动态与跨区域滞后信息，并通过门控融合形成互补修正。

去除BCMS后，联合MAE和联合RMSE分别增加0.005\%和0.064\%。相较于两类专家，BCMS在全测试集平均指标上提供了较小的增益。结合第\ref{sec:scenario-bcms}节的背景差异分组分析，其作用进一步体现为随交通背景变化的历史消息补充。因此，消融结果与情景分析从整体性能和条件响应两个层面共同支持了模型设计：SERR组织自身演化与有向时滞证据，BCMS在此基础上提供背景条件化的进一步修正。

% !TeX root = ../main.tex
\begin{table}[!htbp]
  \centering
  \caption{成都T5数据集上的消融测试结果}
  \label{tab:serr-expert-ablation}
  \begingroup
  \linespread{1}\fontsize{10}{12}\selectfont
  \setlength{\tabcolsep}{5pt}
  \renewcommand{\arraystretch}{1.05}
  \begin{tabular*}{\linewidth}{@{\extracolsep{\fill}}lrrrr@{}}
    \toprule
    模型配置 & 联合MAE$\downarrow$ & 相对变化 & 联合RMSE$\downarrow$ & 相对变化 \\
    \midrule
    CMSRN & \textbf{0.708537} & --- & \textbf{1.258507} & --- \\
    w/o Self Expert (A1) & 0.711270 & +0.386\% & 1.262301 & +0.301\% \\
    w/o Delay Expert (A2) & 0.708852 & +0.044\% & 1.263068 & +0.362\% \\
    w/o BCMS (A3) & 0.708570 & +0.005\% & 1.259318 & +0.064\% \\
    \bottomrule
  \end{tabular*}
  \par\vspace{3pt}
  \begin{minipage}{\linewidth}
  \footnotesize
  注：各变体均独立重新训练。联合指标为流入、流出AVG12指标的等权平均；相对变化以完整CMSRN的对应指标为参照，正值表示误差增加。加粗表示列内最低误差。
  \end{minipage}
  \endgroup
\end{table}

\FloatBarrier
